\section{Numerical verification of the product rate}\label{sec:numerics}

The one new random-matrix constant in the paper --- the product least-singular-value rate
$(k+1)H(\alpha)/\dF$ of \Cref{thm:product} --- differs from the previously conjectured $2H(\alpha)/\dF$
starting at $k=2$, so it is worth checking directly. We do so by exhaustive minimization over all
support tuples at small $N$, where the exponent can be read from a single matrix, and by the cascade
construction at slightly larger $N$, where the achieving pattern can be built but exhaustive search is
infeasible.

\subsection{Setup}
For $N\in\{8,10\}$, density $\alpha=1/2$ (so $n=N/2$), and real Gaussian entries $N(0,1/N)$, we compute
the single-layer minimum $M_N^{(1)}=\min_{|S_0|=|S_1|=n}\smin(A_1[S_1,S_0])$ and the depth-two minimum
$M_N^{(2)}=\min_{S_0,S_1,S_2}\smin(A_2[S_2,S_1]A_1[S_1,S_0])$ over all support triples, reporting the
rates $\rho_k=-N^{-1}\log M_N^{(k)}$. For $N=8$ the triple minimization is over $70^3$ tuples and is
done exactly; for $N=10$ it is over $252^3$ tuples, so there we report the greedy cascade (the
two-sided-optimal first factor followed by the one-sided-optimal second factor), which is the
achieving construction of \Cref{app:product} and upper-bounds the exhaustive minimum. The discriminating
statistic is the ratio $\rho_2/\rho_1$: \Cref{thm:product} predicts $3H/2H=1.5$, the earlier conjecture
$2H/2H=1$.

\subsection{Results}
The measured rates are as follows (means over $8$ trials at $N=8$, $4$ at $N=10$; $2H(1/2)=1.386$,
$3H(1/2)=2.079$).

\begin{center}
\renewcommand{\arraystretch}{1.2}
\begin{tabular}{@{}lccc@{}}
\toprule
 & $\rho_1$ (exhaustive) & $\rho_2$ (exhaustive) & $\rho_2/\rho_1$ \\
\midrule
$N=8$  & $1.375$ & $2.087$ & $1.53$ \\
$N=10$ & $1.579$ & $2.338$ (cascade) & $1.50$ (cascade) \\
\bottomrule
\end{tabular}
\end{center}

At $N=8$, where the depth-two minimum is exact, $\rho_2=2.087$ sits on the predicted $3H=2.079$, and the
ratio $\rho_2/\rho_1=1.53$ is at the value $1.5$ that \Cref{thm:product} requires, cleanly separated from
the conjectured $1.0$. The greedy cascade tracks the exhaustive minimum from above (its rate is a lower
estimate of the true rate, and lies within finite-$N$ error), confirming that the joint-support optimization of \Cref{app:product} is the mechanism realizing the
minimum. At $N=10$ the cascade ratio is $1.50$. The single-layer rate $\rho_1$ fluctuates around
$2H=1.386$ at these small sizes (across seeds, means of $1.42$, $1.31$, $1.36$ at $N=8,10,12$), the
$o(N)$ correction to a first-passage exponent being of either sign; the ratio $\rho_2/\rho_1$, which
cancels the leading correction, is the robust quantity and confirms the linear-in-$k$ growth of the
exponent.

\subsection{The anisotropic cascade step}
The first moment of the cascade construction (\Cref{prop:aniso-mean}) predicts that the least singular
value of an anisotropic row ensemble $M$ with covariance $\Sigma$ of least eigenvalue $\lambda_{\min}$,
minimized over row supports, scales as $\sqrt{n\lambda_{\min}}$ times the isotropic rate. We check the
$\sqrt{\lambda_{\min}}$ prefactor directly. Taking $\Sigma=\diag(1,\dots,1,\eps^2)$ and computing the
exact minimum over row supports at $N=10$, $n=5$, the minimum is proportional to $\eps$ once $\eps$ is
small enough that the least eigenvalue dominates: the ratio $\min/(\eps\cdot\min|_{\eps=1})$ saturates to
a constant for $\eps\le 0.1$, and the log--log slope of $\min$ against $\eps$ is $1.0$--$1.15$ there,
matching the predicted linear dependence. This is the regime the cascade uses, where
$\lambda_{\min}=\smin(Q)^2$ is exponentially small.

The code is \texttt{numerics/verify\_product\_rate.py} and \texttt{verify\_anisotropic\_onesided.py};
both use only \texttt{numpy}. The verification refutes the product conjecture of \cite{shmaloProducts}
and supports \Cref{thm:product} and the first-moment reduction of \Cref{prop:aniso-mean}.
